% GENERATED FILE. Edit canonical lessons, then run npm run generate:books.
% canonical-source-hash: fnv1a64:696a8a209b4ce52a
% canonical-lessons: JA-C252-okite, JA-C252-tabete-kara, JA-C252-sorekara, JA-C252-soshite, JA-C252-tsukihi, JA-W252-tsuki-kanji, JA-R252-one-thing-then-the-next

\chapter{One Thing, Then the Next}
\label{ch:ja-one-thing-then-the-next}
\hlchaptermodality{252}

\begin{chapteropening}
\textbf{By the end of this chapter:} \emph{I can tell what I did in order, joining actions with the -te form (okite, kao o aratte, dekakemashita), say what came after what (tabete kara; sore kara; soshite), and recognise \ja{月} in \ja{月日}.}

Complete JA-R252-one-thing-then-the-next: tell a morning in order, one action after another.
\end{chapteropening}

\section[\texorpdfstring{\ja{おきて} / \ja{いって} (okite / itte)}{おきて / いって (okite / itte)}]{\ja{おきて} / \ja{いって} (okite / itte) — one action, then the next}

\label{lesson:JA-C252-okite}



\begin{quote}
Say, to a friend, \textquotedblleft{}I got up\textquotedblright{} and \textquotedblleft{}I went\textquotedblright{}. (\emph{okita}, \emph{itta}.) Then say \textquotedblleft{}I washed my face\textquotedblright{}, politely. (\emph{kao o araimashita}.)
\end{quote}

\subsection*{You'll want to know: \ja{た} becomes \ja{て}}
\begin{quote}
\textbf{\ja{おきて、かおを} \ja{あらいました。}}
\end{quote}

\begin{quote}
\emph{okite, kao o araimashita.} — \textquotedblleft{}I got up and washed my face.\textquotedblright{}
\end{quote}

\textbf{\ja{おきて}} (\emph{okite}) is \textbf{\ja{おきた}} (\emph{okita}) with \textbf{\ja{え}} in place of \textbf{\ja{あ}}: \textbf{\ja{た}} becomes \textbf{\ja{て}}. A sentence can hold several actions this way, in the order they happened, and the last verb says when: \textbf{\ja{ました}} puts all of them in the past.

\begin{quote}
\textbf{\ja{あさ} \ja{おきて、かおを} \ja{あらいました。}}
\end{quote}

\begin{quote}
\emph{asa okite, kao o araimashita.} — \textquotedblleft{}In the morning I got up and washed my face.\textquotedblright{}
\end{quote}

The other verbs that drop \textbf{\ja{る}} do the same, and so do \textbf{\ja{する}} and \textbf{\ja{くる}}: \textbf{\ja{たべた}} gives \textbf{\ja{たべて}} (\emph{tabete}), \textbf{\ja{した}} gives \textbf{\ja{して}} (\emph{shite}), \textbf{\ja{きた}} gives \textbf{\ja{きて}} (\emph{kite}).

\begin{quote}
\textbf{\ja{ばんごはんを} \ja{たべて、ねました。}}
\end{quote}

\begin{quote}
\emph{bangohan o tabete, nemashita.} — \textquotedblleft{}I ate dinner and went to bed.\textquotedblright{}
\end{quote}

You know this form already from a request: \textbf{\ja{いって} \ja{ください}} (\emph{itte kudasai}), \textquotedblleft{}please say it\textquotedblright{}.

\subsection*{You'll want to know: \ja{って}, \ja{んで}, \ja{いて}, \ja{して}}
The past by the last kana gives the \textbf{-\ja{て}} form by the last kana, with the same swap: \textbf{\ja{だ}} becomes \textbf{\ja{で}}.

\begin{itemize}
\raggedright
  \item \textbf{\ja{いった}} $\to$ \textbf{\ja{いって}} (\emph{itte}), \textbf{\ja{かった}} $\to$ \textbf{\ja{かって}} (\emph{katte}), \textbf{\ja{あらった}} $\to$ \textbf{\ja{あらって}} (\emph{aratte})
  \item \textbf{\ja{のんだ}} $\to$ \textbf{\ja{のんで}} (\emph{nonde}), \textbf{\ja{あそんだ}} $\to$ \textbf{\ja{あそんで}} (\emph{asonde})
  \item \textbf{\ja{かいた}} $\to$ \textbf{\ja{かいて}} (\emph{kaite}), \textbf{\ja{およいだ}} $\to$ \textbf{\ja{およいで}} (\emph{oyoide}), \textbf{\ja{はなした}} $\to$ \textbf{\ja{はなして}} (\emph{hanashite})
\end{itemize}

\begin{quote}
\textbf{\ja{かおを} \ja{あらって、でかけました。}}
\end{quote}

\begin{quote}
\emph{kao o aratte, dekakemashita.} — \textquotedblleft{}I washed my face and went out.\textquotedblright{}
\end{quote}

\begin{quote}
\textbf{\ja{えきに} \ja{いって、きっぷを} \ja{かいました。}}
\end{quote}

\begin{quote}
\emph{eki ni itte, kippu o kaimashita.} — \textquotedblleft{}I went to the station and bought a ticket.\textquotedblright{}
\end{quote}

\subsection*{Your turn}
Say these aloud:

\begin{itemize}
\raggedright
  \item \textquotedblleft{}I drank tea and went out\textquotedblright{} — \emph{ocha o nonde, dekakemashita}
  \item \textquotedblleft{}I drank coffee and went out\textquotedblright{} — \emph{kōhī o nonde, dekakemashita}
  \item \textquotedblleft{}I drank water and went out\textquotedblright{} — \emph{mizu o nonde, dekakemashita}
\end{itemize}

\begin{tcolorbox}[breakable,colback=teal!4,colframe=teal!35!black,title={Before you move on}]
The \textbf{-\ja{て}} form of \textbf{\ja{おきる}}? (\textbf{\ja{おきて}}.) Of \textbf{\ja{する}}? (\textbf{\ja{して}}.) Of \textbf{\ja{いく}}? (\textbf{\ja{いって}}.) Of \textbf{\ja{のむ}}? (\textbf{\ja{のんで}}.) Which verb says when it happened? (The last one.)
\end{tcolorbox}
\section[\texorpdfstring{\ja{たべてから} (tabete kara)}{たべてから (tabete kara)}]{\ja{たべてから} (tabete kara) — after doing}

\label{lesson:JA-C252-tabete-kara}



\begin{quote}
Say \textquotedblleft{}I ate breakfast and went out\textquotedblright{}. (\emph{asagohan o tabete, dekakemashita}.) Now make the order matter: first breakfast, and only then out.
\end{quote}

\subsection*{You'll want to know: -\ja{てから}}
\begin{quote}
\textbf{\ja{あさごはんを} \ja{たべてから、でかけました。}}
\end{quote}

\begin{quote}
\emph{asagohan o tabete kara, dekakemashita.} — \textquotedblleft{}After eating breakfast, I went out.\textquotedblright{}
\end{quote}

\textbf{\ja{から}} (\emph{kara}) after the \textbf{-\ja{て}} form makes \textquotedblleft{}after\textquotedblright{}: the first action was over before the second began. \textbf{-\ja{て}} alone lines actions up; \textbf{-\ja{てから}} says one waited for the other.

\begin{quote}
\textbf{\ja{てがみを} \ja{かいてから、ねました。}}
\end{quote}

\begin{quote}
\emph{tegami o kaite kara, nemashita.} — \textquotedblleft{}After writing a letter, I went to bed.\textquotedblright{}
\end{quote}

\subsection*{You'll want to know: plans and habits}
With the non-past, \textbf{-\ja{てから}} gives the order of a plan or a habit:

\begin{quote}
\textbf{\ja{はを} \ja{みがいてから、ねます。}}
\end{quote}

\begin{quote}
\emph{ha o migaite kara, nemasu.} — \textquotedblleft{}I brush my teeth before I go to bed.\textquotedblright{}
\end{quote}

\begin{quote}
\textbf{\ja{しゅくだいを} \ja{してから、あそびます。}}
\end{quote}

\begin{quote}
\emph{shukudai o shite kara, asobimasu.} — \textquotedblleft{}I will play after doing my homework.\textquotedblright{}
\end{quote}

Two more words for \textquotedblleft{}and then\textquotedblright{}, which start a new sentence, come next.

\subsection*{Your turn}
Say these aloud:

\begin{itemize}
\raggedright
  \item \textquotedblleft{}After eating dinner, I went to bed\textquotedblright{} — \emph{bangohan o tabete kara, nemashita}
  \item \textquotedblleft{}After eating lunch, I went to bed\textquotedblright{} — \emph{hirugohan o tabete kara, nemashita}
  \item \textquotedblleft{}After eating breakfast, I went to bed\textquotedblright{} — \emph{asagohan o tabete kara, nemashita}
\end{itemize}

\begin{tcolorbox}[breakable,colback=teal!4,colframe=teal!35!black,title={Before you move on}]
\textquotedblleft{}After eating\textquotedblright{}? (\textbf{\ja{たべてから}}.) \textquotedblleft{}After writing\textquotedblright{}? (\textbf{\ja{かいてから}}.) What does \textbf{\ja{から}} add to \textbf{-\ja{て}}? (The first action is over before the next begins.)
\end{tcolorbox}
\section[\texorpdfstring{\ja{それから} (sorekara)}{それから (sorekara)}]{\ja{それから} (sorekara) — and then; after that}

\label{lesson:JA-C252-sorekara}



\begin{quote}
Before the new word: say \textquotedblleft{}after eating breakfast, I went out\textquotedblright{}. (\emph{asagohan o tabete kara, dekakemashita}.) Then say \textquotedblleft{}that\textquotedblright{}, a thing near the listener. (\emph{sore}.)
\end{quote}

\subsection*{You'll want to know: \ja{それから}}
\textbf{\ja{それから}} — \emph{sorekara} — \textquotedblleft{}and then\textquotedblright{}, \textquotedblleft{}after that\textquotedblright{}. It opens the next sentence, and \textbf{\ja{、}} follows it.

\begin{quote}
\textbf{\ja{あさごはんを} \ja{たべました。}}
\end{quote}

\begin{quote}
\emph{asagohan o tabemashita.} — \textquotedblleft{}I ate breakfast.\textquotedblright{}
\end{quote}

\begin{quote}
\textbf{\ja{それから、でかけました。}}
\end{quote}

\begin{quote}
\emph{sorekara, dekakemashita.} — \textquotedblleft{}After that, I went out.\textquotedblright{}
\end{quote}

\begin{cousinweb}
\textbf{\ja{それから}} (\emph{sorekara}) is a native Japanese word (wago), inherited rather than borrowed from Chinese: \textbf{\ja{それ}} (\emph{sore}), \textquotedblleft{}that\textquotedblright{}, and \textbf{\ja{から}} (\emph{kara}), \textquotedblleft{}from\textquotedblright{}, so \textquotedblleft{}from that point on\textquotedblright{}. JMdict lists it as an expression, \textquotedblleft{}and then; after that\textquotedblright{} (JMdict).
\end{cousinweb}

\subsection*{Your turn}
Say these aloud:

\begin{itemize}
\raggedright
  \item \emph{sorekara}
  \item \textquotedblleft{}After that, I went out\textquotedblright{} — \emph{sorekara, dekakemashita}
  \item \textquotedblleft{}After that, I went home\textquotedblright{} — \emph{sorekara, kaerimashita}
\end{itemize}

\begin{tcolorbox}[breakable,colback=teal!4,colframe=teal!35!black,title={Before you move on}]
\textquotedblleft{}And then\textquotedblright{}, at the start of a sentence? (\textbf{\ja{それから}}, \emph{sorekara}.) Which mark follows it? (\textbf{\ja{、}}.)
\end{tcolorbox}
\section[\texorpdfstring{\ja{そして} (soshite)}{そして (soshite)}]{\ja{そして} (soshite) — and; joining two sentences}

\label{lesson:JA-C252-soshite}



\begin{quote}
Before the new word: say \textquotedblleft{}after that, I went out\textquotedblright{}. (\emph{sorekara, dekakemashita}.)
\end{quote}

\subsection*{You'll want to know: \ja{そして}}
\textbf{\ja{そして}} — \emph{soshite} — \textquotedblleft{}and\textquotedblright{}, at the start of a sentence. \textbf{\ja{それから}} says what came next in time; \textbf{\ja{そして}} adds the next thing.

\begin{quote}
\textbf{\ja{うみに} \ja{いきました。}}
\end{quote}

\begin{quote}
\emph{umi ni ikimashita.} — \textquotedblleft{}I went to the sea.\textquotedblright{}
\end{quote}

\begin{quote}
\textbf{\ja{そして、およぎました。}}
\end{quote}

\begin{quote}
\emph{soshite, oyogimashita.} — \textquotedblleft{}And I swam.\textquotedblright{}
\end{quote}

\begin{cousinweb}
\textbf{\ja{そして}} (\emph{soshite}) is a native Japanese word (wago), inherited rather than borrowed from Chinese. It is usually explained as \textbf{\ja{そ}} (\emph{so}), \textquotedblleft{}that way\textquotedblright{}, the same \emph{so-} as in \textbf{\ja{それ}}, with \textbf{\ja{して}} (\emph{shite}), the \textbf{-\ja{て}} form of \textbf{\ja{する}}. JMdict marks it as usually written in kana (JMdict).
\end{cousinweb}

\subsection*{Your turn}
Say these aloud:

\begin{itemize}
\raggedright
  \item \emph{soshite}
  \item \textquotedblleft{}And I swam\textquotedblright{} — \emph{soshite, oyogimashita}
  \item \textquotedblleft{}And I played\textquotedblright{} — \emph{soshite, asobimashita}
\end{itemize}

\begin{tcolorbox}[breakable,colback=teal!4,colframe=teal!35!black,title={Before you move on}]
\textquotedblleft{}And\textquotedblright{}, opening a sentence? (\textbf{\ja{そして}}, \emph{soshite}.) \textquotedblleft{}And then\textquotedblright{}, in time? (\textbf{\ja{それから}}, \emph{sorekara}.)
\end{tcolorbox}
\section[\texorpdfstring{\ja{月日} (tsukihi)}{月日 (tsukihi)}]{\ja{月日} (tsukihi) — the passing of time}

\label{lesson:JA-C252-tsukihi}



\begin{quote}
Before the new word: say \textquotedblleft{}the moon\textquotedblright{}. (\emph{tsuki}.) Then say \textquotedblleft{}fast\textquotedblright{}. (\emph{hayai}.)
\end{quote}

\subsection*{You'll want to know: \ja{月日}}
\textbf{\ja{月日}} — \emph{tsukihi} — \textquotedblleft{}the passing of time\textquotedblright{}: months and days, gone by.

\begin{quote}
\textbf{\ja{月日}\ja{は} \ja{はやいです。}}
\end{quote}

\begin{quote}
\emph{tsukihi wa hayai desu.} — \textquotedblleft{}Time flies.\textquotedblright{}
\end{quote}

Two beats, then two more: \emph{tsu–ki–hi}. The second sign is \textbf{\ja{日}}, the kanji you write. The first is new: \textbf{\ja{月}}, read \textbf{\ja{つき}} (\emph{tsuki}) here, \textquotedblleft{}the moon\textquotedblright{}, the word you know. The next lesson writes it.

\begin{cousinweb}
\textbf{\ja{月日}} (\emph{tsukihi}) is a native Japanese word (wago), inherited rather than borrowed from Chinese: \textbf{\ja{つき}} (\emph{tsuki}), \textquotedblleft{}the moon, a month\textquotedblright{}, and \textbf{\ja{ひ}} (\emph{hi}), \textquotedblleft{}a day\textquotedblright{}, each written with its kanji and read the native way, as a kun'yomi. JMdict glosses the whole as \textquotedblleft{}time; years; days\textquotedblright{}, and KANJIDIC2 gives \ja{つき} for the first kanji and \ja{ひ} for the second (JMdict, KANJIDIC2).
\end{cousinweb}

\subsection*{Your turn}
\begin{itemize}
\raggedright
  \item \emph{Say it:} \emph{tsukihi}
  \item \emph{Say it:} \textquotedblleft{}Time flies\textquotedblright{} — \emph{tsukihi wa hayai desu}
  \item \emph{Recall:} say the moon, then a day, then the two together
\end{itemize}

\begin{tcolorbox}[breakable,colback=teal!4,colframe=teal!35!black,title={Before you move on}]
\textquotedblleft{}The passing of time\textquotedblright{}? (\textbf{\ja{月日}}, \emph{tsukihi}.) Which sign do you already write? (\textbf{\ja{日}}.) How is the new one read here? (\textbf{\ja{つき}}, \emph{tsuki}.)
\end{tcolorbox}
\section[\texorpdfstring{\ja{月} (tsuki)}{月 (tsuki)}]{\ja{月} — the moon, four strokes}

\label{lesson:JA-W252-tsuki-kanji}



\begin{quote}
Say \emph{tsukihi}, the passing of time. Its second sign is \textbf{\ja{日}}, which you write. Today the first.
\end{quote}

\subsection*{Script — the shape}
\begin{quote}
\textbf{\ja{月}} — \emph{tsuki}
\end{quote}

Its parts:

\begin{itemize}
\raggedright
  \item a tall frame, open at the bottom
  \item two short bars inside
\end{itemize}

Four strokes, three lifts:

\begin{enumerate}
\raggedright
  \item draw the left side from the top down, sweeping a little to the left at the foot
  \item lift, draw the top from left to right, then without lifting turn down the right side and hook left at its foot
  \item lift, draw the upper inside bar from left to right
  \item lift, draw the lower inside bar from left to right
\end{enumerate}

Like \textbf{\ja{日}} it is taller than it is wide and has bars inside, but it is open at the bottom, and its first stroke sweeps away to the left. Here it is read \textbf{\ja{つき}} (\emph{tsuki}), as in \textbf{\ja{月日}}. Kanji have more than one reading; this lesson asks only for the one in the word you know.

\subsection*{Your turn}
\begin{enumerate}
\raggedright
  \item Find \textbf{\ja{月}} in \textbf{\ja{月日}} and say \emph{tsukihi}.
  \item Copy \textbf{\ja{月}} three times, saying \emph{tsuki} at each one.
  \item Find it once more, in \textbf{\ja{月日}\ja{は} \ja{はやいです}} (\emph{tsukihi wa hayai desu}), and say the line.
\end{enumerate}

\begin{tcolorbox}[breakable,colback=teal!4,colframe=teal!35!black,title={Before you move on}]
Which way does the first stroke sweep at its foot? (A little to the \textbf{left}.)

Source: \href{https://github.com/KanjiVG/kanjivg/blob/master/kanji/06708.svg}{KanjiVG stroke-order data}: KanjiVG, kanji/06708.svg, kanji \ja{月} (U+6708), 4 stroke paths kvg:06708-s1 to s4 (Ulrich Apel and contributors, CC BY-SA 3.0). Only the order and direction are taken from it.
\end{tcolorbox}
\section[(review)]{Review — one thing, then the next}

\label{lesson:JA-R252-one-thing-then-the-next}



\begin{quote}
Say \textquotedblleft{}I got up and washed my face\textquotedblright{}, then \textquotedblleft{}after that, I went out\textquotedblright{}. (\emph{okite, kao o araimashita}; \emph{sorekara, dekakemashita}.)
\end{quote}

\subsection*{Your turn}
Say these aloud:

\begin{itemize}
\raggedright
  \item In the morning I got up, and washed my face.
  \item After eating breakfast, I went out.
  \item I went to the sea. And I swam.
  \item I wrote a letter. After that, I went to bed.
  \item Time flies.
\end{itemize}

\begin{tcolorbox}[breakable,colback=teal!4,colframe=teal!35!black,title={Before you move on}]
The \textbf{-\ja{て}} form of \textbf{\ja{かう}}? (\textbf{\ja{かって}}.) Of \textbf{\ja{よむ}}? (\textbf{\ja{よんで}}.) \textquotedblleft{}After writing\textquotedblright{}? (\textbf{\ja{かいてから}}.) \textquotedblleft{}After that\textquotedblright{}, opening a sentence? (\textbf{\ja{それから}}.) \textquotedblleft{}And\textquotedblright{}, opening a sentence? (\textbf{\ja{そして}}.) \textquotedblleft{}Time flies\textquotedblright{}? (\textbf{\ja{月日}\ja{は} \ja{はやいです}}.)
\end{tcolorbox}
